The domain source is Azuonye's study of the Nwagu Aneke Igbo script, which describes the script's origins, its syllabic character, logographs, limitations, and literacy potential \cite{Azuonye1992}. Azuonye is therefore the anchor for the source audit, but the paper does not itself validate PAGC's later 27x8, E6, or compression claims.

The compression side of the archive uses Byte Pair Encoding as a falsification tool. BPE-style subword segmentation is well established in neural machine translation \cite{Sennrich2016}; here it is used only as a local test of the k=27 optimum claim.

The infrastructure side is related to scientific knowledge graphs and claim verification. The Open Research Knowledge Graph represents scholarly knowledge in structured form \cite{Jaradeh2019}, and recent scientific claim-verification benchmarks frame support checking as an explicit task \cite{Ho2026}. PAGC differs from those settings because the central failure is not retrieval alone; it is a foundation-count drift between an archived visual source, derived theoretical notation, and paper claims.

The workflow also draws on automated research loops such as autoresearch \cite{Karpathy2026}, but with a different risk model: rather than optimizing an experiment metric, the loop blocks claims whose source layer, evidence locator, or contradiction status is unclear. The packaging constraints follow current arXiv and ACM guidance \cite{ArxivTeXLive2026,ACMSubmissions2026}.
